problem:  
Let \((M_i^n,g_i)\) be a sequence of closed Riemannian manifolds with \(\operatorname{Ric}_{M_i}\ge -(n-1)\), \(\operatorname{Diam}(M_i)\le D\), and \( \operatorname{Vol}(M_i)\to V>0\).  Assume \((M_i,g_i)\) converge in the measured Gromov–Hausdorff sense to \((X,d_X,\mathcal{H}^n)\); by Cheeger–Colding and Mondino–Naber, \(X\) is an \(\mathrm{RCD}(-(n-1),n)\) space.  

For each \(i\), fix a smooth isoperimetric minimizer \(\Omega_i\subset M_i\) for some fixed reference volume fraction \(v\in (0,1)\), i.e.  
\[
\operatorname{Vol}_{M_i}(\Omega_i)=v\,\operatorname{Vol}(M_i),\qquad
\operatorname{Area}_{g_i}(\partial\Omega_i)=I_{M_i}(v). 
\]
By compactness of the sequence \((M_i,g_i)\), the characteristic functions \(\chi_{\Omega_i}\) form a sequence of BV functions with uniformly bounded variation.  Hence (after subsequence extraction) they may converge in \(L^1\) to a set of finite perimeter \(\Omega_\infty\subset X\).  

**Conjectural problem:**  
Does the limit perimeter measure of \(\Omega_i\) detect the singular structure of \(X\)?  More precisely, is it true or false that the perimeter measure \( |D\chi_{\Omega_i}|\) concentrates asymptotically only on the regular set \(\mathcal{R}\subset X\), in the sense that the limiting measure of the singular part vanishes:
\[
\lim_{i\to\infty}\frac{|D\chi_{\Omega_i}|(f_i^{-1}(\mathcal{S}))}{|D\chi_{\Omega_i}|(M_i)} = 0,
\]
where \(\mathcal{S}=X\setminus\mathcal{R}\) and \(f_i:M_i\to X\) are almost isometries realizing the measured Gromov–Hausdorff convergence?  

Equivalently: in Ricci-limit spaces, do genuine isoperimetric regions in the approximating manifolds “see” only the smooth part of the limit, or can they exhibit nontrivial area concentration along the singular strata?  

Why is it a "good" problem:  
This question isolates a subtle and genuinely unresolved aspect of Colding’s volume convergence philosophy—it asks whether the *fine variational geometry* of isoperimetric hypersurfaces remains smooth in the limit or interacts significantly with metric singularities.  Existing BV compactness and lower-semicontinuity results in the RCD setting do not decide this issue: they guarantee convergence of perimeters but not the *distribution* of perimeter measure between regular and singular parts.  Resolving this would deepen our understanding of how curvature lower bounds control the interface between smooth and singular regions of Ricci-limit spaces.  

A positive answer would demonstrate that lower Ricci curvature bounds preclude area concentration on singular strata, thus unifying geometric measure theory with the stratified analysis of Ricci-limit spaces and advancing structural classification beyond mere metric or measure convergence.  A counterexample would reveal new analytic phenomena in singular metric geometry.  The problem is expressible in a single line yet probes the core of the Cheeger–Colding–Naber theory, exemplifying the ideal of simplicity masking profound depth.